本书的目标是让你成为更优秀的程序员。

我写这本书，是为了把我从数十年的专业软件开发与教学生涯中所学到的东西传承下去。我成年后的人生一直在硅谷学习、生活、工作和教书。我曾在成熟的计算机公司和几家初创公司担任高级工程职位。我曾在 IBM 研究院和劳伦斯利弗莫尔国家实验室（Lawrence Livermore National Laboratory）开发过先进软件，还曾是 NASA 和 JPL 的高级科学家。我曾在多所大学为本科生和研究生教授软件开发，其中包括圣何塞州立大学（San Jose State University），它为硅谷输送的工程师比任何其他学校都多。

与学生和初级程序员一起工作让我明白，在坏习惯养成之前就实践良好的软件设计是很重要的。作为学生，我们无意中学到了“跑通即完成”：只要被布置的程序能成功运行，就算完成了！把它交上去之后，我们可能再也不用见到它，因此诸如可维护性这样的良好设计概念就变得无关紧要。要想在软件开发者这条职业道路上取得成功，我们就必须摒弃这种心态。

我非常清楚按时、按预算完成应用所面临的压力。因此，在本书中，我也教授一种迭代式、增量式的软件开发方法。如果由于临近截止日期而无法完成最后一轮迭代、交付完整的产品，我们希望能由倒数第二轮（或倒数第三轮）迭代的成果产出一个 MVP（最小可行产品）。但愿产品还有下一个版本，可以用来清理设计问题并添加更多功能。

\subsection{经验是最好的老师}

除了学习本书中的示例之外，成为一流程序员的关键是什么？那就是练习、练习、再练习！我希望“经验是最好的老师”这句话既适用于作为软件开发者的你，也适用于作为教师和写作者的我。
